\chapter{見えない「場」の支配と情報の伝播 --- 電磁気学とニューラルネットワークの波}

\section*{本章の学習目標}
\begin{itemize}
\item ニュートンの「遠隔作用」の謎と、物理学における「空間」の捉え方の限界を知る。
\item ファラデーの直感とマクスウェルの方程式による「場（Field）」の発見の歴史を学ぶ。
\item 「ポテンシャル」「div（湧き出し）」「rot（渦）」といったベクトル解析の幾何学的な意味を理解する。
\item 「光の正体」が電磁波であることを解き明かした、物理学史上の奇跡の瞬間を体感する。
\item 全結合（遠隔作用）からCNN（近接作用・受容野）への、AIの視覚革命の軌跡を理解する。
\item 現代AIの心臓部「Self-Attention（自己注意機構）」が作り出す「動的な意味の場」の美しさを解き明かす。
\end{itemize}

\section{ニュートンの限界：オカルトと呼ばれた「遠隔作用」}

第1章で触れた通り、アイザック・ニュートンが打ち立てた「万有引力の法則」は、地上でリンゴが落ちる法則と、天空で月が地球の周りを回る法則が「全く同じ数式」で説明できることを証明し、人類の宇宙観を永遠に変えました。

しかし、この完璧に見えるニュートン力学には、当時から多くの学者（そしてニュートン自身でさえも）が深く思い悩む「致命的で気味が悪い弱点」が一つだけ隠されていました。

\subsection{何もない真空を「瞬時」に伝わる力}

地球と月との間には、約38万キロメートルもの「何もない真空の空間」が広がっています。それにもかかわらず、地球と月は互いに目に見えない糸で結ばれているかのように、強烈な引力で引き合っています。

もし今、神様が太陽を突然消し去ったとしたら、ニュートンの数式に従えば、約1億5000万キロメートル離れた地球の軌道は「その瞬間に（時間差ゼロで）」変化し、宇宙の果てへ向かって飛んでいってしまうことになります。

「なぜ、間には何もないのに、はるか遠く離れた物体同士の力が『瞬時』に伝わるのか？」

触れてもいないのに力が伝わるこの奇妙な現象を、物理学では\textbf{「遠隔作用（Action at a distance）」}と呼びます。

当時の科学者たちにとって、物質的な媒介なしに力が伝わるというのは、魔法やオカルトと何ら変わらない不気味な考え方でした。ニュートン自身も友人への手紙の中で、「遠く離れた物体が、何もない空間を飛び越えて互いに作用を及ぼすなど、常識を備えた人間には信じがたい不条理だ」と書き残しているほどです。

\subsection{デカルトの「渦動説」との激しい対立}

事実、ニュートンの万有引力が発表された当時、ヨーロッパ大陸の知識人たち（ライプニッツやホイヘンスなど）はこれを猛烈に批判しました。

当時の科学界の主流は、フランスの哲学者ルネ・デカルトが提唱した\textbf{「渦動説（かどうせつ）」でした。デカルトは「宇宙には真空など存在せず、目に見えない微細な物質がぎっしり詰まっていて、それが巨大な渦を巻いているから惑星は押し流されて回っているのだ」と主張していました。これは、歯車が噛み合って動く時計のように、物質同士が直接「接触」して力を伝える「近接作用」}という非常に常識的でわかりやすい考え方です。

それに引き換え、ニュートンの「真空を飛び越えて目に見えない力が引っ張り合う」という主張は、中世の魔法使いが呪文を唱えて遠くのものを動かすような、時代遅れのオカルトにしか見えなかったのです。「ニュートンは科学を中世の暗黒時代に引き戻そうとしている！」とまで非難されました。

\subsection{錬金術師ニュートンの素顔と「我、仮説を立てず」}

なぜニュートンは、同時代の科学者が嫌悪したような「目に見えない神秘的な力」を数式化できたのでしょうか？

実は、ニュートンは近代科学の父であると同時に、生涯の膨大な時間を\textbf{「錬金術」}や聖書の暗号解読に費やした、ヨーロッパ最後の魔術師とも呼べる人物でした。水銀や硫黄を煮詰め、「賢者の石」や「隠された霊的な力」を本気で信じて研究していた彼だからこそ、何もない空間を伝わる不気味な「引力」という概念を、直感的に受け入れることができたとも言われています。

批判に対し、ニュートンは著書の中で\textbf{「我、仮説を立てず（Hypotheses non fingo）」}という有名な言葉を残しました。

「引力が『なぜ』発生するのか、その根本的なメカニズム（仮説）は私には分からないし、捏造するつもりもない。しかし、引力がこの数式通りに『どう動くか』は完璧に証明したのだから、今はそれで十分ではないか」と開き直ったのです。

力はどのようにして空間を渡っていくのか？ 重力の正体とは何なのか？ この巨大なミステリー（WhatとHowは分かるがWhyが分からない、まさに現代AIのブラックボックスと同じ状態）が解き明かされるまでには、ニュートンの死後、約150年もの歳月を待たねばなりませんでした。

\section{ファラデーの直感と魔法：「場（Field）」の発見}

19世紀、この不条理な「遠隔作用」の呪縛を打ち破ったのは、エリートの数学者や大学教授ではなく、正規の数学教育を一切受けていない、イギリスの貧しい製本所の見習い職人でした。彼の名をマイケル・ファラデーと言います。

\subsection{製本職人の見習いと「数式なき天才」}

ファラデーは小学校すらろくに卒業しておらず、微積分はおろか、簡単な代数（数式）すら理解できませんでした。しかし、この「数学の欠如」こそが、逆に彼に物理学史上最大のパラダイムシフトを起こさせる武器となりました。

当時のエリート学者たちは、ニュートンの重力方程式に倣い、電気や磁気の力を「数式上の遠隔作用」として計算することに終始していました。しかし数式が読めないファラデーは、数式に頼る代わりに、自らの手で行う実験と、圧倒的な\textbf{「視覚的な想像力（幾何学的な直感）」}だけで自然界の秘密を解き明かそうとしたのです。

\subsection{砂鉄が描く見えない線}

ファラデーは、磁石や電気の実験に異常なまでの情熱を持っていました。

紙の上に砂鉄を撒き、その下に棒磁石を置くと、砂鉄は磁石のN極からS極へと向かう美しい曲線状の模様を描き出します。誰もが子供の頃に見たことのある光景です。

当時の物理学者たちは、これを「磁石のN極とS極が、砂鉄の粒一つ一つに対して直接『遠隔作用』で引力を及ぼしている結果に過ぎない」と計算していました。

しかし、直感の天才であったファラデーには、全く違う世界が見えていました。彼は、磁石が遠くの砂鉄を直接引っ張っているのではなく、\textbf{「磁石が存在することで、その周りにある『何もない空間そのもの』に目に見えない緊張や歪みが生じているのだ」}と考えたのです。

\subsection{「近接作用」と電磁誘導：空間を伝わるドミノ倒し}

ファラデーは、この空間に張り巡らされた見えない力の線を\textbf{「力線（Lines of force）」、そして力が存在する空間そのものを「場（Field）」}と名付けました。

この「場」の概念は、物理学の世界観を根本からひっくり返しました。

力は、遠くへ一瞬で魔法のように飛んでいく（遠隔作用）のではありません。磁石がまず自分の「すぐ隣の場」を歪ませ、その歪みがさらに隣へ隣へと、まるで池に投げ入れた小石の波紋や、ドミノ倒しのように伝わっていくのです。そして、場を伝わってきた歪みが最終的に砂鉄の置かれた場所に到達したとき、初めて砂鉄に力が働きます。

これを\textbf{「近接作用（Local action）」}と呼びます。何もないと思われていた真空の空間は、実は力を伝える見えないゴムシートのような「場」で満たされていたのです。

さらにファラデーは1831年、この「場」の概念を決定づける大発見をします。コイル（導線の束）の中で磁石を動かすと、コイルに電流が流れる\textbf{「電磁誘導」}の法則です。これは、「磁石が動くことで空間の『場』の歪みが時間とともに変化し、その波紋が空間を伝わって導線の電子を押し流した」という、動的な近接作用の決定的な証拠でした。

\subsection{光の「偏光」の発見と横波の証明}

ファラデーの直感がさらに飛躍する前に、当時の物理学界で起きていたもう一つの「光の謎解き」のドラマに触れておきましょう。

17世紀以来、光が「粒子」なのか「波」なのかは激しい論争の的でしたが、19世紀初頭には「光は波である」という説が優勢になっていました。しかし、光が「音のように進行方向と同じ向きに揺れる波（縦波）」なのか、「水面のように進行方向に対して垂直に揺れる波（横波）」なのかは決着がついていませんでした。

この謎を偶然の発見で打ち破ったのが、フランスの技術将校エティエンヌ＝ルイ・マリュスです。1808年のある日、彼はパリのリュクサンブール宮殿の窓ガラスに反射する夕日を、特殊な結晶（方解石）を通して眺めていました。すると、結晶を回す角度によって、夕日の光が明るくなったり、完全に消えて真っ暗になったりすることを発見したのです。

マリュスは、\textbf{「光は単なる波ではなく、縦や横といった特定の『向き』を持って振動している」ことに気づきました。これを「偏光（へんこう：Polarization）」}と呼びます。

その後、同じくフランスの物理学者オーギュスタン・ジャン・フレネルらがこの偏光現象を数学的に徹底的に分析し、「光は進行方向に対して垂直に振動する『横波』でなければ、この現象は絶対に説明できない」ことを完璧に証明しました。

\subsection{偏光と磁気の融合：ファラデー効果と「光線振動説」}

晩年、ファラデーはこの「光の偏光」という不思議な現象と、自らが研究してきた「磁場」の力とを頭の中で結びつけようとしました。「もし空間が本当に『場』という見えないシートでできているのなら、磁石で作った場の歪みの中を光（波）が通るとき、光の波の揺れ方に何らかの影響が出るはずだ」。

1845年、ファラデーは強力な電磁石の間にガラスを置き、そこに「特定の向き（偏光面）を持った光」を通す実験を行いました。すると予想通り、強力な磁場を通った光は、その偏光面（揺れる向き）が磁場の力によってクルリと回転したのです。これを\textbf{「ファラデー効果」}と呼びます。全く別物だと思われていた「光」と「磁気」が、空間の「場」を介して直接相互作用することを、人類で初めて実験で証明した瞬間でした。

翌1846年、興奮冷めやらぬファラデーは『光線振動説（Thoughts on Ray-vibrations）』という短い論文を発表し、歴史的な大予言を行います。

「宇宙空間に張り巡らされた電気や磁気の力線がブルブルと震え、その波紋が伝わっていく現象こそが『光』の正体なのではないか？」

しかし、このファラデーの革新的な予言は、当時のアカデミアからは冷たくあしらわれました。彼が高度な数式を使えず、直感的な言葉だけで論文を書いたため、「物理学ではなく、素人のロマンチックな空想に過ぎない」とバカにされたのです。

\section{場の数学：ポテンシャル、湧き出し（div）、渦（rot）}

ファラデーが直感で描いた「見えない力線」の幻視を、完璧な数学の言語へと翻訳するためには、空間全体の状態を記述するための全く新しい数学が必要でした。それが、マクスウェルやその他の天才たちが磨き上げた\textbf{「ベクトル解析」}という数学分野です。

彼らは、目に見えない電気や磁気の「場」を、まるで\textbf{「空間を流れる巨大な川の水流」}のようにイメージしました。この水流の振る舞いを記述するために、彼らは3つの極めて直感的で強力な概念（数学的オペレーター）を導入しました。

\subsection{ポテンシャル（高さ）とグラディエント（傾き：\(\mathrm{grad}\)または\(\nabla\)）}

第1章や第2章で、ボールが高いところから低いところへ転がる「エネルギーの山谷」を見ました。「場」の世界でも、空間の各点に見えないエネルギーの「高さ」が設定されています。これを\textbf{「ポテンシャル（例えば電位や電圧）」}と呼びます。

水が高いところから低いところへ流れるように、このポテンシャルの「傾き（グラディエント）」が急であればあるほど、そこに置かれた物質（電荷など）は強い「力」を受けて転がり落ちます。私たちがスマホの充電に使う「電圧」とは、まさにこの電気的なポテンシャルの高低差のことです。これは、AIが「勾配降下法（グラディエント・ディセント）」で損失の谷底へ向かうときにも現れる、\textbf{「傾きが次の変化を決める」}という共通の数学的構造です。

\subsection{ダイバージェンス（湧き出し：\(\mathrm{div}\)または\(\nabla\cdot\)）}

空間のある一点から、水（力線）が外に向かってどれくらい湧き出しているか、あるいは吸い込まれているかを表す指標です。お風呂の蛇口（プラスの湧き出し）や、底の排水溝（マイナスの吸い込み）を想像してください。物理学では、電荷（プラスやマイナス）がまさにこの「電場の湧き出し口や吸い込み口」として機能します。

\subsection{ローテーション（渦：\(\mathrm{rot}\)あるいは\(\mathrm{curl}\)、\(\nabla\times\)）}

空間のある一点で、水の流れがどれくらい「渦」を巻いているかを表す指標です。川の中に小さな水車を置いたとき、左右の水流のスピードの違いによって、水車がどれくらい勢いよくクルクルと回るかを示します。

マクスウェルは、この「湧き出し（div）」と「渦（rot）」という2つの幾何学的な視点を使うことで、それまで無関係だと思われていた数々の実験結果が、実は美しいパズルのピースであることを見抜いたのです。

\section{巨星たちのバケツリレー：マクスウェル方程式に至る4つのピース}

孤独なファラデーの直感の真価を見抜き、その「空想」を完璧なベクトル解析の言葉に翻訳したのが、スコットランドの若き天才理論物理学者、ジェームズ・クラーク・マクスウェルでした。しかし、彼がゼロからすべての方程式を創り出したわけではありません。電磁気学の完成は、18世紀から19世紀にかけて活躍した巨星たちによる、数式と実験のバケツリレーの賜物でした。

マクスウェルが統合した「電気と磁気に関する4つのピース（法則）」を、実際の数式（微分形）とともに見ていきましょう。（\(\mathbf{E}\)は電場、\(\mathbf{B}\)は磁場を表し、\(\nabla\cdot\)（div：湧き出し）、\(\nabla\times\)（rot：渦）の記号を使います）

\subsection{ピース1：電荷が「電場」の湧き出し口になる（ガウスの法則）}

18世紀末、フランスのクーロンは「電荷（プラスとマイナス）」の間に働く力を定式化しました。これを「場」の視点から洗練させたのが、数学の王様カール・フリードリヒ・ガウスです。

\[
\nabla\cdot\mathbf{E}=\frac{\rho}{\varepsilon_0}
\]

数式の意味：空間にある電荷\(\rho\)からは、電場\(\mathbf{E}\)がウニのトゲやホースの水のように四方八方へ「湧き出して（div）」いる。

\subsection{ピース2：磁気には「単独の湧き出し口」がない（磁場に対するガウスの法則）}

電気にはプラス単独、マイナス単独の湧き出し口が存在しますが、磁石をいくら細かく割っても、必ずN極とS極がペアで現れます。

\[
\nabla\cdot\mathbf{B}=0
\]

数式の意味：磁場\(\mathbf{B}\)には、N極単独のような「湧き出し口（div）」は存在せず、湧き出し量は常にゼロである。磁力線は必ずループして閉じる。

\subsection{ピース3：電流が「磁場の渦」を作る（アンペールの法則）}

1820年、デンマークのエルステッドが「導線に電流を流すと方位磁針が振れる」ことを発見し、フランスのアンドレ＝マリ・アンペールがこれを数学的に定式化しました。

\[
\nabla\times\mathbf{B}=\mu_0\mathbf{J}
\]

数式の意味：電流\(\mathbf{J}\)が流れると、その周囲にぐるぐると「渦を巻くような磁場（\(\nabla\times\mathbf{B}\)：rot）」が発生する。

\subsection{ピース4：磁場の変化が「電場の渦」を作る（ファラデーの電磁誘導の法則）}

アンペールらの発見を知ったファラデーは、「電気が磁気を作るなら、磁気で電気を作れるはずだ」と考え、コイルの中で磁石を動かすことで電流が流れる「電磁誘導」を発見しました。

\[
\nabla\times\mathbf{E}=-\frac{\partial\mathbf{B}}{\partial t}
\]

数式の意味：磁場が時間とともに変化する（\(\frac{\partial \mathbf{B}}{\partial t}\)）と、空間に「渦を巻くような電場（rot）」が発生する。

\section{マクスウェルの「変位電流」と光の正体}

マクスウェルは、これら4つの数式を並べて眺めていたとき、アンペールの法則（ピース3）に致命的な数学的矛盾があることに気づきました。途切れた回路（コンデンサなど）を流れる電流を計算しようとすると、数式の辻褄が合わなくなってしまったのです。

\subsection{対称性への美意識と「最後のピース」}

ここでマクスウェルは、物理学者としての並外れた美的直感を発揮します。

ファラデーの法則（ピース4）によれば、「磁場の変化」は空間に「電場の渦（rot）」を生み出します。ならば自然界の美しい対称性から考えて、逆に\textbf{「電場の変化」もまた空間に「磁場の渦（rot）」を生み出すはずだ}、と考えたのです。

彼は実験的な証拠がまだ何もないにもかかわらず、アンペールの法則の右辺に「電場が時間とともに変化する割合（\(\frac{\partial \mathbf{E}}{\partial t}\)）」という全く新しい架空の電流（変位電流）の項を強制的に付け足しました。

\[
\nabla\times\mathbf{B}=\mu_0\mathbf{J}+\mu_0\varepsilon_0\frac{\partial\mathbf{E}}{\partial t}
\]

この「変位電流」というたった一つの項が加わったことで、4つの方程式の間の矛盾は消え去り、物理学の至宝\textbf{「マクスウェルの方程式」}がここに完成したのです。

この項の意味は、「導線の中を電荷が流れていなくても、電場そのものが時間的に変化していれば、それは磁場を作る原因として働く」ということです。つまり、マクスウェルは物質の流れだけでなく、\textbf{場の変化そのものが場を生む}という、電磁気学の自己増殖的な構造を方程式に刻み込んだのです。

\subsection{電磁波の「波動方程式」の導出}

マクスウェルの計算は、衝撃的な事実を浮き彫りにしました。

空間に電荷も電流もない真空中（\(\rho=0,\mathbf{J}=0\)）を考えてみましょう。このとき、第3式と第4式は次のような完全に美しい対称な形になります。

\[
\nabla\times\mathbf{E}=-\frac{\partial\mathbf{B}}{\partial t}
\]

\[
\nabla\times\mathbf{B}=\mu_0\varepsilon_0\frac{\partial\mathbf{E}}{\partial t}
\]

マクスウェルは、この2つの式を数学的に組み合わせました。第3式の両辺についてさらに「渦（\(\nabla\times\)）」の計算を行い、そこに第4式を代入します。ベクトル解析の公式（\(\nabla\times(\nabla\times\mathbf{E})=\nabla(\nabla\cdot\mathbf{E})-\nabla^2\mathbf{E}\)）と第1式（\(\nabla\cdot\mathbf{E}=0\)）を巧みに使うと、最終的に次のようなたった一つの方程式にたどり着きます。

\[
\nabla^2\mathbf{E}=\mu_0\varepsilon_0\frac{\partial^2\mathbf{E}}{\partial t^2}
\]

（※磁場\(\mathbf{B}\)についても全く同じ形の方程式が導かれます）

物理学において、この形の方程式（空間の2階微分が、時間の2階微分に比例する）は\textbf{「波動方程式」と呼ばれ、水面やロープを伝わる「波」の動きを完全に表すものです。 つまり、「電場が変化する」→「その隣に磁場の渦（rot）が生まれる」→「磁場が変化する」→「その隣に電場の渦（rot）が生まれる」という相互作用の連鎖が、空間を伝わる「波（電磁波）」}として存在することが、純粋な数学的計算から自動的に証明されたのです。

\subsection{物理学史上の奇跡：光速度は「定数」だけで決まる}

さらに劇的だったのは、マクスウェルがこの「電磁波の波及するスピード」を計算したときでした。

一般的に波動方程式\(\nabla^2 u=\frac{1}{v^2}\frac{\partial^2u}{\partial t^2}\)において、右辺の時間微分の前にある係数\(\frac{1}{v^2}\)の逆数の平方根（\(v\)）が「波の伝わるスピード」を表します。マクスウェルの導き出した電磁波の波動方程式にこれを当てはめると、電磁波のスピード（\(c\)）は次のような極めてシンプルな数式で計算されます。

\[
c=\frac{1}{\sqrt{\varepsilon_0\mu_0}}
\]

驚くべきことに、この式には「波の形」や「周波数」などは一切含まれていません。電磁波のスピードは、\textbf{「真空の誘電率（\(\varepsilon_0\)：電気の通しにくさ）」と「真空の透磁率（\(\mu_0\)：磁気の通しにくさ）」}という、宇宙空間の根本的な性質を表す「たった2つの定数」だけで完全に決まってしまうのです。

方程式にこの2つの定数の値を入れて計算を終えたマクスウェルは、紙の上に現れた数字を見て息を呑みました。

計算によって弾き出された電磁波のスピードは「秒速約30万キロメートル」。

それは、当時の別の実験ですでに測定されていた\textbf{「光のスピード」と、小数点以下の桁に至るまで完全に一致していた}のです。

「光とは、空間の『場』を伝播する電磁波の一種だったのだ！」

ファラデーの予言した「光線振動説」は、決して素人の空想などではなく、宇宙の真理そのものでした。電気の力、磁石の力、そして私たちが目で見ている光。これら全く別物だと思われていた3つの現象が、マクスウェル方程式という一つの美しい数式の下で完璧に統合された、人類の知性の歴史における奇跡の瞬間でした。（そしてこの「光の速度は誰から見ても一定である」というマクスウェルの数式の特徴が、数十年後にアインシュタインの相対性理論を生み出す決定的な火種となります）。

\section{AIにおけるパラダイムシフト：全結合から「受容野」へ}

空間に対する捉え方が、ニュートンの「遠隔作用」から、ファラデーとマクスウェルによる「場を通じた近接作用」へと進化した物理学の歴史。

実はこれと全く同じ、そして極めて劇的な歴史的パラダイムシフトが、現代のAI（人工知能）が「視覚（画像認識能力）」を獲得する進化の過程でも起きています。

\subsection{「全結合」の罠：遠隔作用の限界と空間構造の破壊}

1980年代から90年代にかけて研究されていた初期のニューラルネットワーク（多層パーセプトロン）で画像を認識させようとしたとき、研究者たちは致命的な問題に直面していました。

当時のAIは、画像のすべてのピクセルを入力として横一列に並べ、次の層のすべてのニューロンと「総当たり（全結合）」で線を結んで計算していました（これを全結合層：Fully Connected Layerと呼びます）。

これは物理学で言えば、まさにニュートンの\textbf{「遠隔作用」}と全く同じ状態です。画面の右上のピクセルと左下のピクセルが、物理的な「距離」という概念を完全に無視して、直接関係性を計算しようとするのです。さらに、ピクセルの数が増えれば増えるほど、すべての点と点を結ぶためのパラメータ（線の数）が天文学的な数に爆発してしまい、計算が破綻してしまいます。

何より致命的だったのは、この方法では「犬の耳の形」や「車のタイヤの丸み」といった、「隣り合うピクセル同士の空間的な繋がり（構造）」が完全に破壊されてしまうことでした。画像を「バラバラの点の集まり」としてしか認識できない全結合のAIは、画像に写る犬の位置が右に少しズレただけで、全く犬だと認識できなくなるという絶望的な脆弱性を抱えていました。

\subsection{CNNと「受容野（Receptive Field）」：近接作用の誕生}

この遠隔作用の限界を打破し、2012年のディープラーニング革命（AIの視覚革命）の直接的な引き金となったのが、生物の視覚野からヒントを得て作られた\textbf{「CNN（畳み込みニューラルネットワーク）」}という技術です。

CNNは、画面全体をバラバラに一度に見ることをやめました。代わりに、3×3ピクセル程度の小さな四角い\textbf{「フィルター（虫眼鏡のようなもの）」を用意し、それを画像の上で少しずつスライドさせながら計算を行っていきます。 つまり、遠く離れたピクセルをいきなり結びつけるのではなく、「すぐ隣り合うピクセルとの関係性だけを計算し、その波紋を次の層へ渡していく」というアプローチを取ったのです。これはまさに、ファラデーが提唱した「空間の歪みが隣へ隣へと伝わっていく」という「場（近接作用）」の概念そのものへの転換}でした。

このフィルターが見つめている局所的な空間範囲のことを、AIの用語で\textbf{「受容野（Receptive Field）」}と呼びます。

第1層の受容野は3×3の狭い「場」しか見ておらず、単なる線や角（エッジ）しか認識できません。しかし、層が深くなるにつれて、近接作用のドミノ倒しのように情報が統合され、第2層、第3層と進むにつれてより広い「場」を見るようになります。そして最終的には、犬の顔や車の全体像といった大域的な意味を理解できるようになるのです。

\subsection{並進対称性と「宇宙のどこでも同じ法則」}

CNNが物理学と見事にシンクロしているもう一つの点が、\textbf{「重み共有（Weight Sharing）」}という仕組みです。

CNNの3×3のフィルターは、画像の左上を見る時も、右下を見る時も、全く同じ計算ルール（重み）を使い回します。

これは、物理学における最も重要で美しい公理の一つである\textbf{「空間の並進対称性（宇宙のどこに行っても物理法則は同じである）」}と完全に一致しています。犬の耳は、写真の左上に写っていても、右下に写っていても「犬の耳」として同じように認識されなければなりません。CNNは、フィルターの重みを共有することで、この物理法則が持つ「空間的な対称性」をAIのネットワーク構造そのものに組み込んだのです。これにより、対象物がどこにズレても認識できる堅牢性と、計算パラメータの劇的な削減を実現しました。

\subsection{偏微分とエッジ抽出：数学的な双子}

さらに驚くべきことに、マクスウェル方程式とCNNは、その心臓部の\textbf{「計算の数学的構造」までもが完全に同じ}です。

前節で見たマクスウェル方程式には、「\(\mathrm{rot}\)（\(\nabla\times\)：渦）」や「\(\mathrm{div}\)（\(\nabla\cdot\)：湧き出し）」といった記号が登場しました。これらはベクトル解析における\textbf{「偏微分（空間の微小な変化を捉える計算）」です。コンピュータ上でこの偏微分を計算するときは、ある点の値から「すぐ隣の点の値」を引き算するという「空間の差分」}をとります。

一方、CNNの第1層がやっていることは、画像の中から「エッジ（輪郭線）」を見つけ出すことです。エッジとは「画像の明るさが急激に変化している場所」のことです。CNNのフィルター（例えばソーベルフィルタ）は、あるピクセルの明るさから「すぐ隣のピクセルの明るさ」を引き算して、その変化の度合いを計算します。

お気づきでしょうか。CNNのフィルターが行っている「畳み込み積分（Convolution）」によるエッジ抽出は、マクスウェル方程式を解くための「偏微分（差分）」と全く同じ幾何学的な操作なのです。

「遠く離れたものをいきなり結びつけるのではなく、すぐ隣同士の変化（微分）を計算し、その影響を次へと伝えていく」。

19世紀の物理学者たちが宇宙空間の秘密を解き明かすために発明した「場の数学」は、21世紀のコンピュータの中で、AIが世界を視覚的に理解するための「究極のアルゴリズム」として完璧に蘇り、人間を超える圧倒的な認識能力を手に入れたのです。

\section{万物と繋がる力：Self-Attentionと「動的な場」}

AIにおける「場」の概念は、CNNの成功にとどまらず、現在世界を席巻しているChatGPTなどの大規模言語モデル（LLM）の心臓部において、究極の形へと昇華されています。それが、2017年にGoogleの研究者たちが発表したTransformerアーキテクチャの中核をなす\textbf{「Self-Attention（自己注意機構）」}です。

\subsection{意味の空間を満たす「Attentionの場」}

文章というものは、単語が一列に並んだ一つの「空間」だと考えることができます。

例えば、「The animal didn't cross the street because it was too tired.（その動物は道を渡らなかった、なぜなら『それ』は疲れすぎていたからだ）」という文章があったとします。

この時、「it（それ）」という単語は、ストリート（道）のことでしょうか？ それともアニマル（動物）のことでしょうか？ 人間は文脈から「疲れるのは動物だ」と一瞬で理解しますが、これをコンピュータに計算させるのは至難の業でした。

第4.6節で見たCNNのような「隣り合うピクセルだけを見る固定されたフィルター（近接作用）」は、このように遠く離れた場所にある単語同士の「意味の繋がり」をうまく掴むことができません。

そこで2017年、Googleの研究者たちは物理学の「場」の概念を究極の形に進化させました。それがSelf-Attentionです。

彼らは、入力されたすべての単語を、意味の空間（多次元のベクトル空間）に配置された\textbf{「電荷（電気を帯びた粒）」のように扱いました。そして、ある単語（it）が、文章全体のすべての単語に対して「互いにどれくらい強く関連し合っているか（引力を及ぼし合っているか）」のスコア（Attention Weight）を一斉に計算}させたのです。

\subsection{Query、Key、Value：共鳴する電磁波}

この相互作用を計算するために、Self-Attentionは各単語に3つの役割（ベクトル）を持たせます。

Query（クエリ：問いかけ）：「私は『疲れた』という性質を持つ名詞を探しています」という、自分が空間に放つ電磁波のようなサイン。

Key（キー：鍵）：「私は『動物』という属性を持っています」という、周囲に自らの性質を知らせるアンテナ。

Value（バリュー：中身）：その単語が持っている実質的な「意味」のデータそのもの。

空間の中で、ある単語のQuery（発信）と、別の単語のKey（受信）の波長がピタリと共鳴したとき（数学的には2つのベクトルの内積が大きくなったとき）、両者の間に強力な\textbf{「意味の引力」}が生じます。

上記の例文では、「it」が発したQueryの波と、「animal」が発していたKeyの波が最も強く共鳴します。すると「it」は「street」よりも「animal」に強い引力を感じ取り、「animal」の持つ意味（Value）を自分の中に強く吸収するのです。

\subsection{究極の「動的な場」が知能を生む}

Self-Attentionの凄まじさは、CNNのように「あらかじめ決まった3×3の範囲の場だけを見る」のではなく、\textbf{「入力された文章の文脈に応じて、単語同士が相互作用し、毎回全く新しい『意味の力場（相互作用のネットワーク）』をリアルタイムに形成する」}という点にあります。

これは一見すると、ニュートンの「遠隔作用（距離を無視してすべての点が直接結びつく全結合）」に逆戻りしたように見えるかもしれません。しかし本質は全く違います。距離を無視しているのではなく、\textbf{「物理的な距離ではなく『意味的な距離』に応じた力場を、その瞬間の文脈に合わせて動的に創り出している」}のです。

これはまさに、空間に配置された無数の電荷が互いに電磁力を及ぼし合い、瞬時に引力と斥力を計算して、空間全体に複雑でダイナミックな「電磁場」を形成するマクスウェル方程式の世界そのものです。

数千億のパラメータを持つChatGPTのニューラルネットワークの深海では、一つ一つの単語（トークン）という粒が互いに共鳴し合い、複雑な「Attentionの波」を何十層にもわたって伝播させていきます。層を重ねるごとに、単語は周囲の単語の意味を吸い込んで進化し、最終的に「文脈全体」という巨大な意味の大海原を完璧に計算し尽くしているのです。

「遠隔作用」から「近接作用の場」へ、そして「動的に相互作用する場」へ。

物理学が「何もない空間」に隠された力の伝播ルールを解き明かしたように、AIもまた、データの背後に潜む「場を通じた意味の絡み合い」を数学的にシミュレートすることで、私たち人間と対話できるほどの圧倒的な知性を獲得するに至ったのです。

\section*{【第4章 章末ディスカッション課題】}

ニュートンが信じられなかった「遠隔作用」と、私たちがインターネットで一瞬にして地球の裏側と通信できることは、何が違うのでしょうか？ 私たちの社会は「場」によってどのように支えられていると思いますか？

人間の脳が世界を認識するとき、それはCNNのように「局所的な細かいパーツ（場）を組み合わせて全体を作っている」のでしょうか？ それとも、Attentionのように「最初から全体を同時に関連付けて」意味を理解しているのでしょうか？ 人間の直感はどちらに近いと思いますか？


\section*{参考文献（学術誌）}
\begin{enumerate}[leftmargin=2.4em]
\item Maxwell, J. C. ``A dynamical theory of the electromagnetic field.'' \textit{Philosophical Transactions of the Royal Society of London}, 155, 459--512, 1865. DOI: \url{https://doi.org/10.1098/rstl.1865.0008}.
\item Hodgkin, A. L. and Huxley, A. F. ``A quantitative description of membrane current and its application to conduction and excitation in nerve.'' \textit{The Journal of Physiology}, 117(4), 500--544, 1952. DOI: \url{https://doi.org/10.1113/jphysiol.1952.sp004764}.
\item Scarselli, F., Gori, M., Tsoi, A. C., Hagenbuchner, M., and Monfardini, G. ``The graph neural network model.'' \textit{IEEE Transactions on Neural Networks}, 20(1), 61--80, 2009. DOI: \url{https://doi.org/10.1109/TNN.2008.2005605}.
\item Bronstein, M. M., Bruna, J., LeCun, Y., Szlam, A., and Vandergheynst, P. ``Geometric deep learning: going beyond Euclidean data.'' \textit{IEEE Signal Processing Magazine}, 34(4), 18--42, 2017. DOI: \url{https://doi.org/10.1109/MSP.2017.2693418}.
\end{enumerate}


